%% 2.1.4 Kerberos: obter o TGT
\begin{tikzpicture}
  \node[vm, text width=3.20cm, font=\footnotesize] (u) at (-5.75,0) {Ubuntu\\(\texttt{kinit})};
  \node[vm, text width=3.20cm, font=\footnotesize] (d) at (0.00,0) {DNS\\(\texttt{SR2-SRV})};
  \node[vm, text width=3.20cm, font=\footnotesize] (k) at (5.75,0) {KDC\\(\texttt{SR2-SRV}:88)};
  \draw[vida] (u.south) -- (-5.75,-7.53);
  \draw[vida] (d.south) -- (0.00,-7.53);
  \draw[vida] (k.south) -- (5.75,-7.53);
  \draw[seta] (-5.75,-1.60) -- (0.00,-1.60);
  \node[rot, anchor=south] at (-2.88,-1.56) {Onde está o KDC? (SRV \texttt{\_kerberos.\_tcp})};
  \draw[setad] (0.00,-2.43) -- (-5.75,-2.43);
  \node[rot, anchor=south] at (-2.88,-2.39) {\texttt{SR2-SRV.isec.local}};
  \draw[seta] (-5.75,-3.26) -- (5.75,-3.26);
  \node[rot, anchor=south] at (0.00,-3.22) {AS-REQ \texttt{aluno.teste@ISEC.LOCAL}};
  \draw[setad] (5.75,-4.09) -- (-5.75,-4.09);
  \node[rot, anchor=south] at (0.00,-4.05) {Pré-autenticação necessária};
  \draw[seta] (-5.75,-4.92) -- (5.75,-4.92);
  \node[rot, anchor=south] at (0.00,-4.88) {AS-REQ com prova da palavra-passe};
  \draw[setad] (5.75,-5.75) -- (-5.75,-5.75);
  \node[rot, anchor=south] at (0.00,-5.71) {AS-REP com o TGT (\texttt{krbtgt/ISEC.LOCAL})};
  \node[acao, text width=2.9cm] at (-5.75,-6.63) {Guarda o TGT na cache (\texttt{klist})};
\end{tikzpicture}
